\chapter{Flexibility and rigidity of geometric structures}

\begin{definition}[Holonomy deformation equivalence]
\label{thurston-ch05-holonomy-deformation-equivalence}
\group{deformation theory}\source{thurston:page-0097}\notready
For a compact $(G,X)$--manifold $M$, two nearby structures are equivalent when, after deleting small neighborhoods of their boundaries, a near-identity $(G,X)$--homeomorphism identifies them.  The developing map determines a holonomy representation $\pi_1M\to G$.
\end{definition}

\begin{proposition}[Local holonomy parametrization]
\label{thurston-ch05-local-holonomy}
\dcref{Proposition 5.1}\group{deformation theory}\source{thurston:page-0097,thurston:page-0098}\notready
\uses{thurston-ch05-holonomy-deformation-equivalence}
Nearby $(G,X)$--structures on a compact manifold are classified, up to the preceding equivalence, by nearby holonomy representations modulo conjugation by small elements of $G$.
\end{proposition}
\begin{proof}\sketch Represent $M$ by a disk with overlap data on its boundary.  Perturbing the holonomy perturbs the finitely many identifications and hence gives a nearby structure.  Conjugate holonomies are made equal by a small coordinate change; restriction to smaller disks then supplies the required equivalence.\end{proof}

\begin{proposition}[Generic representation dimension]
\label{thurston-ch05-generic-representation-dimension}
\group{deformation theory}\source{thurston:page-0098}\notready
If a presentation has $|\G|$ generators and $|\mathcal R|$ relators, and the relator map has surjective derivative at a representation, its local representation space has dimension $(|\G|-|\mathcal R|)\dim G$.  When conjugation has full-dimensional orbits, the local deformation dimension is $\dim G(|\G|-|\mathcal R|-1)$.
\end{proposition}
\begin{proof}\sketch Embed representations in $G^{\G}$, identify them as the identity fiber of the relator map, apply the implicit function theorem, and quotient by the conjugation orbit.\end{proof}

\begin{proposition}[Three-dimensional algebraic lower bound]
\label{thurston-ch05-three-dimensional-lower-bound}
\dcref{Proposition 5.2.2}\group{deformation theory}\source{thurston:page-0099,thurston:page-0100}\notready
If a hyperbolic $3$--manifold $M$ exists, its small deformation space has dimension at least $6(-\chi(M))$.
\end{proposition}
\begin{proof}\sketch The relator equations define a complex algebraic fiber in a product of copies of $\PSL(2,\C)$; its dimension is at least variables minus equations.  The cell structure gives the stated Euler-characteristic count.\end{proof}

\begin{proposition}[Closed geodesic representatives]
\label{thurston-ch05-closed-geodesic-representatives}
\dcref{Proposition 5.3.1}\group{hyperbolic geometry}\source{thurston:page-0100,thurston:page-0101}\notready
Every nontrivial free homotopy class in a closed hyperbolic manifold contains a unique closed geodesic.
\end{proposition}
\begin{proof}\sketch A nontrivial deck transformation is neither elliptic nor parabolic: ellipticity contradicts freeness and parabolicity would give arbitrarily short closed curves in a compact manifold.  It is therefore hyperbolic, translating along a unique axis; axes project to the same geodesic exactly under conjugacy.\end{proof}

\begin{proposition}[Disjoint geodesic realization]
\label{thurston-ch05-disjoint-geodesic-realization}
\dcref{Proposition 5.3.3}\group{hyperbolic surfaces}\source{thurston:page-0101}\notready
Pairwise disjoint, homotopically distinct essential simple closed curves on a closed hyperbolic surface are represented by pairwise disjoint simple closed geodesics.
\end{proposition}
\begin{proof}\sketch Intersecting lifts have linked distinct endpoints on the circle at infinity; finite-distance homotopies cannot change endpoints, so the intersection cannot be removed.  Applying the same argument to distinct lifts proves simplicity.\end{proof}

\begin{proposition}[Pair of pants coordinates]
\label{thurston-ch05-pants-coordinates}
\dcref{Proposition 5.3.4}\group{Teichmuller theory}\source{thurston:page-0102}\notready
The Teichmuller space of a hyperbolic pair of pants is $\R^3$, with coordinates $(\log l_1,\log l_2,\log l_3)$ given by its positive geodesic boundary lengths.
\end{proposition}
\begin{proof}\sketch The three perpendicular seams split the pants into two right-angled hexagons.  Alternating side lengths determine each hexagon, and reflection identifies the halves, so arbitrary positive boundary lengths determine one and only one structure.\end{proof}

\begin{theorem}[Fenchel--Nielsen coordinates]
\label{thurston-ch05-fenchel-nielsen}
\dcref{Theorem 5.3.5}\group{Teichmuller theory}\source{thurston:page-0101,thurston:page-0102,thurston:page-0103,thurston:page-0104}\notready
For a closed genus-$g$ surface, $\mathcal T(M)\cong\R^{6g-6}$, with coordinates $(\log l_i,\tau_i)_{i=1}^{3g-3}$ consisting of lengths and real twist parameters along a pants decomposition.  A full twist changes the marked point even when it preserves the unmarked isometry class.
\end{theorem}
\begin{proof}\sketch Cut along $3g-3$ geodesics into $2g-2$ pants, assign the three length coordinates on each piece, and glue matching boundaries with one real twist per interior curve.  The Euler characteristic count gives the number of curves.\end{proof}

\begin{theorem}[Fricke and torus coordinates]
\label{thurston-ch05-fricke-torus-coordinates}
\dcref{Theorem 5.3.6,Theorem 5.3.7}\group{Teichmuller theory}\source{thurston:page-0104}\notready
For a hyperbolic surface with geodesic boundary, the Fricke space is $\R^{-3\chi(M)}$; boundary curves carry no twist coordinate.  The Teichmuller space of a unit-area flat torus is $\R^2$ with coordinates $(\log l,\tau)$ in a chosen primitive free homotopy class.
\end{theorem}

\begin{proposition}[Jorgensen word symmetries]
\label{thurston-ch05-jorgensen-word-symmetries}
\dcref{Proposition 5.4.1}\group{PSL algebra}\source{thurston:page-0105,thurston:page-0106}\notready
If $a,b\in\PSL(2,\C)$ have no common fixed point at infinity and a word $w(a,b)$ is trivial, then $w(a^{-1},b^{-1})$ is trivial.  If $a,b$ are conjugate, then $w(b,a)$ is also trivial.
\end{proposition}
\begin{proof}\sketch A half-turn about the canonical common perpendicular (with the analogous limiting construction for parabolics) conjugates both generators to their inverses.  For conjugate generators, half-turns about two orthogonal axes exchange them, possibly after inversion.\end{proof}

\begin{corollary}[Two-generator involutions]
\label{thurston-ch05-two-generator-involutions}
\dcref{Geometric Corollary 5.4.2}\group{hyperbolic 3-manifolds}\source{thurston:page-0106}\notready
A complete hyperbolic $3$--manifold whose group is generated by $a,b$ has an order-two isometry inverting both generators.  If $a,b$ are conjugate and have no common fixed point at infinity, a second commuting involution exchanges them.
\end{corollary}

\begin{proposition}[Common powers and commutation]
\label{thurston-ch05-common-powers-commute}
\dcref{Proposition 5.4.3}\group{PSL algebra}\source{thurston:page-0107}\notready
If non-elliptic isometries satisfy $a^n=b^m\ne1$ for nonzero integers $n,m$, then $a$ and $b$ commute.  If $a$ is hyperbolic and $a^k$ is conjugate to $a^l$, then $k=\pm l$.
\end{proposition}

\begin{proposition}[Centralizers in hyperbolic three-manifolds]
\label{thurston-ch05-centralizer-dichotomy}
\dcref{Proposition 5.4.4}\group{hyperbolic 3-manifolds}\source{thurston:page-0107,thurston:page-0108}\notready
If $a,b\in\pi_1M$ commute in a complete hyperbolic $3$--manifold, then either they lie in one infinite cyclic subgroup, or they lie with finite index in the peripheral group of a torus end.
\end{proposition}
\begin{proof}\sketch Hyperbolic commuting elements share an axis and discreteness makes their translations cyclic.  Otherwise they are parabolic; independent translation vectors generate a Euclidean lattice on a horosphere, and the small-displacement abelian subgroup produces the corresponding torus cusp.\end{proof}

\begin{theorem}[Boundary deformation bound]
\label{thurston-ch05-boundary-deformation-bound}
\dcref{Theorem 5.6}\group{deformation theory}\source{thurston:page-0108,thurston:page-0109,thurston:page-0110,thurston:page-0111,thurston:page-0112}\notready
Let $M$ be compact, oriented, hyperbolic, with nontrivial holonomy on each torus boundary component and no global fixed point on $\Sph^2_\infty$.  Then
\[
\dim_\C\Def(M)\ge\sum_i\begin{cases}3\abs{\chi((\partial M)_i)}&\chi<0,\\1&\chi=0,\\0&\chi>0.\end{cases}
\]
\end{theorem}
\begin{proof}\sketch For every torus choose an element whose holonomy together with the torus subgroup has no common fixed point and is nonparabolic.  Hollowing these curves replaces each torus by a genus-two boundary and two trace equations recover the filling relations.  Counting matrix variables and relations gives the boundary Euler-characteristic bound, with three extra complex dimensions before quotienting by conjugacy.\end{proof}

\begin{theorem}[Mostow rigidity]
\label{thurston-ch05-mostow-rigidity}
\dcref{Theorem 5.7.1,Theorem 5.7.2}\group{rigidity}\source{thurston:page-0113}\notready
For $n\ge3$, isomorphic discrete finite-volume subgroups of $\Isom(\HH^n)$ are conjugate.  Equivalently, every fundamental-group isomorphism between complete finite-volume hyperbolic $n$--manifolds is realized by a unique isometry.
\end{theorem}
\begin{proof}\sketch Lift a homotopy equivalence to a pseudo-isometry of $\HH^n$.  It sends each geodesic a bounded distance from a unique geodesic, inducing a bijection of the spheres at infinity.  The boundary correspondence is continuous and quasiconformal; rigidity of this boundary map forces it to arise from an ambient isometry, which descends uniquely.\end{proof}

\begin{corollary}[Finite outer automorphism group]
\label{thurston-ch05-finite-outer-automorphisms}
\dcref{Corollary 5.7.4}\group{rigidity}\source{thurston:page-0113,thurston:page-0114}\notready
For a complete finite-volume hyperbolic $n$--manifold with $n\ge3$, $\operatorname{Out}(\pi_1M)$ is finite and canonically isomorphic to $\Isom(M)$.
\end{corollary}

\begin{lemma}[Dehn filling from deformed cusp holonomy]
\label{thurston-ch05-dehn-filling-local}
\dcref{Lemma 5.8.1}\group{Dehn surgery}\source{thurston:page-0114,thurston:page-0115}\notready
A sufficiently small deformation of the standard hyperbolic structure on $T^2\times[0,1]$ extends over a solid torus with coefficient $d=(a,b)$, determined up to sign by
\[
a\widetilde H(\alpha)+b\widetilde H(\beta)=\pm2\pi i,
\qquad \widetilde H(\gamma)=\log\lambda_\gamma,
\]
where $\lambda_\gamma+\lambda_\gamma^{-1}$ is the matrix trace.
\end{lemma}
\begin{proof}\sketch Nonparabolic perturbations of the commuting cusp holonomies share a distant axis.  The truncated cusp lifts to the universal cover of $\HH^3$ minus that axis, and metric completion fills its core; the longitudinal and rotational holonomy along the axis gives the displayed coefficient equation.\end{proof}

\begin{theorem}[Hyperbolic Dehn surgery]
\label{thurston-ch05-hyperbolic-dehn-surgery}
\dcref{Theorem 5.8.2}\group{Dehn surgery}\source{thurston:page-0115,thurston:page-0116,thurston:page-0117,thurston:page-0118}\notready
If a finite-ended cusped hyperbolic $3$--manifold $M=M_{\infty,\ldots,\infty}$ exists, then all surgery parameters in a neighborhood of $(\infty,\ldots,\infty)$ produce hyperbolic structures on $M_{d_1,\ldots,d_k}$.  In particular, after excluding finitely many integral slopes on each boundary torus, the remaining primitive surgeries are hyperbolic.
\end{theorem}
\begin{proof}\sketch The boundary deformation bound supplies one complex parameter per cusp.  The trace map is locally open at the unique complete structure by Mostow rigidity.  Near the parabolic locus, the two logarithmic holonomies are asymptotic with ratio determined by the cusp lattice, so the filling equation maps this open trace neighborhood to an open neighborhood of infinity in each slope space.\end{proof}

\begin{lemma}[Pseudo-isometric geodesic correspondence]
\label{thurston-ch05-pseudo-isometric-geodesics}
\dcref{Lemma 5.9.1,Proposition 5.9.2,Corollary 5.9.3}\group{rigidity}\source{thurston:page-0118,thurston:page-0119,thurston:page-0120}\notready
Lifts of homotopy equivalences between closed hyperbolic manifolds can be chosen as pseudo-isometries.  The image of each geodesic lies a bounded distance from a unique geodesic, and this correspondence induces a bijection of the spheres at infinity.
\end{lemma}
\begin{proof}\sketch Compactness bounds the homotopies to the identity, giving the lower pseudo-isometry estimate from the upper Lipschitz estimate.  Hyperbolic projection contracts paths outside a fixed tubular neighborhood, forcing long image geodesics to fellow-travel a limiting line.  Parallelism of lines gives the boundary bijection.\end{proof}

\begin{proposition}[Thick-thin abelian subgroup theorem]
\label{thurston-ch05-thick-thin-abelian-subgroup}
\dcref{Theorem 5.10.1}\group{thick-thin geometry}\source{thurston:page-0124,thurston:page-0125,thurston:page-0126}\notready
There is $\eps>0$ such that, for every discrete group $\Gamma<\Isom(\HH^n)$ and every $x$, the subgroup generated by elements moving $x$ at most $\eps$ has an abelian subgroup of finite index.
\end{proposition}
\begin{proof}\sketch The commutator estimate makes the lower central series of a sufficiently small discrete subgroup terminate, hence makes it nilpotent; for hyperbolic isometries such a subgroup is abelian.  A derivative-pigeonhole argument bounds coset representatives by words of uniformly bounded length, proving finite index.\end{proof}

\begin{corollary}[Thin-part components in dimension three]
\label{thurston-ch05-thin-part-components}
\dcref{Corollary 5.10.2}\group{thick-thin geometry}\source{thurston:page-0127}\notready
For sufficiently small $\eps$, every component of the thin part of a complete oriented hyperbolic $3$--manifold is either a horoball modulo $\Z$ or $\Z\oplus\Z$, or a tubular neighborhood $B_r(g)$ of a geodesic modulo $\Z$; $r=0$ is allowed.
\end{corollary}
\begin{proof}\sketch Small-displacement parabolics generate concentric horoballs, while hyperbolics generate concentric solid cylinders.  Overlap forces commutation by the preceding theorem, so each component has one of these two forms.\end{proof}

\begin{proposition}[Finite-volume cusp decomposition]
\label{thurston-ch05-finite-volume-cusp-decomposition}
\dcref{Proposition 5.11.1}\group{thick-thin geometry}\source{thurston:page-0128}\notready
A complete oriented finite-volume hyperbolic $3$--manifold is the union of a compact submanifold bounded by tori and finitely many rank-two cusps, each a horoball modulo $\Z\oplus\Z$.
\end{proposition}
\begin{proof}\sketch The thick part is compact: infinitely many separated points would yield infinitely many disjoint balls of a fixed positive volume.  Only finitely many thin components meet its compact boundary; absorb the compact ones into the thick part.\end{proof}

\begin{theorem}[Bounded-volume thick-part finiteness]
\label{thurston-ch05-jorgensen-compactness}
\dcref{Theorem 5.11.2}\group{bounded volume}\source{thurston:page-0128,thurston:page-0129,thurston:page-0130}\notready
For every $C>0$, the thick parts $M_{[\eps,\infty)}$ of complete finite-volume hyperbolic $3$--manifolds with volume at most $C$ have only finitely many homeomorphism types.  Consequently, there is a link $L_C\subset S^3$ such that every such manifold is obtained by Dehn surgery along $L_C$, where deleting components is allowed.
\end{theorem}
\begin{proof}\sketch Choose a maximal $\eps/2$-separated set in the thick part.  The disjoint $\eps/4$-balls bound its cardinality in terms of $C$, while the $\eps/2$-balls cover.  Hence only finitely many intersection patterns, and therefore only finitely many thick-part homeomorphism types, occur.  Manifolds with the same thick part differ by Dehn fillings.  Choose one representative for each of the finitely many types, present each representative by surgery on a link in $S^3$, and take the union of these links.\end{proof}

\begin{theorem}[Jorgensen properness theorem]
\label{thurston-ch05-jorgensen-properness}
\dcref{Theorem 5.12.1}\group{bounded volume}\source{thurston:page-0130,thurston:page-0131}\notready
Let $\mathcal H$ be the space of isometry classes of complete finite-volume hyperbolic $3$--manifolds with the geometric topology.  The volume function $v:\mathcal H\to\mathbb R_{>0}$ is proper, so every bounded-volume sequence has a convergent subsequence.  Moreover, for every $C>0$ there are finitely many manifolds $M_1,\ldots,M_k$ of volume at most $C$ such that every other manifold of volume at most $C$ is obtained from one of them by hyperbolic Dehn surgery.
\end{theorem}
\begin{proof}\sketch Use boundedly many balls as above and pass to a subsequence with a fixed gluing pattern.  The gluing isometries range over a compact set, so a further subsequence converges to a hyperbolic thick part.  A geodesic whose length tends to zero has a thin solid-torus neighborhood whose radius tends to infinity; consequently the limiting thick part extends to a complete hyperbolic manifold.  This proves properness, and the description by hyperbolic Dehn surgery is its cusp-filling reformulation.\end{proof}

\begin{proposition}[Solid-torus volume estimate]
\label{thurston-ch05-solid-torus-volume-estimate}
\group{bounded volume}\source{thurston:page-0130,thurston:page-0131}\notready
For a hyperbolic solid torus of core length $l$ and radius $r_0$,
\[
\vol=\pi l\sinh^2r_0,\qquad \Area(\partial)=2\pi l\sinh r_0\cosh r_0,
\]
so $\Area(\partial)/\vol<1/2$, with limiting ratio $1/2$.  Consequently
\[
\vol(M)\le\vol(M_{[\eps,\infty)})+\tfrac12\Area(\partial M_{[\eps,\infty)}).
\]
\end{proposition}
\begin{proof}\sketch Integrate the cylindrical hyperbolic volume density $\sinh r\cosh r$ and the boundary density $l\sinh r_0\cosh r_0$; the ratio follows directly.\end{proof}

\begin{definition}[Injectivity radius]
\label{thurston-ch05-injectivity-radius}
\group{thick-thin geometry}\source{thurston:page-0132,thurston:page-0133}\notready
At $x$ in a Riemannian manifold, $\inj(x)$ is half the infimum of lengths of nontrivial closed loops through $x$, equivalently the supremum of radii on which the exponential map at $x$ is injective.
\end{definition}
